\chapter*{Résumé}
\addcontentsline{toc}{chapter}{Résumé}

Ce mémoire étudie les méthodes numériques à un pas pour la résolution des équations différentielles ordinaires. Après avoir rappelé le cadre théorique du problème de Cauchy (existence et unicité de la solution), nous présentons trois méthodes explicites classiques : la méthode d'Euler, la méthode de Heun et la méthode de Runge--Kutta d'ordre~4.

Nous démontrons qu'une méthode consistante d'ordre~$p$, dont la fonction d'incrément est lipschitzienne, converge à l'ordre~$p$, et nous donnons une estimation explicite de l'erreur globale. Ces résultats sont ensuite confrontés à des expériences numériques : sur une équation test, les ordres observés valent $1$, $2$ et $4$, conformément à la théorie, et le coût nécessaire pour atteindre une erreur de $10^{-6}$ varie de plus de quatre ordres de grandeur d'une méthode à l'autre. Enfin, nous appliquons la méthode de Runge--Kutta au modèle proie-prédateur de Lotka--Volterra et vérifions la conservation de son intégrale première.

\bigskip
\noindent\textbf{Mots-clés :} équations différentielles ordinaires, méthodes à un pas, Runge--Kutta, ordre de convergence, erreur globale.

\clearpage
\begin{otherlanguage}{english}
\chapter*{Abstract}
\addcontentsline{toc}{chapter}{Abstract}

This thesis studies one-step numerical methods for ordinary differential equations. After recalling the theoretical framework of the Cauchy problem (existence and uniqueness of the solution), we present three classical explicit methods: Euler's method, Heun's method and the fourth-order Runge--Kutta method.

We prove that a consistent method of order~$p$ whose increment function is Lipschitz continuous converges with order~$p$, and we give an explicit bound on the global error. These results are then compared with numerical experiments: on a test equation the observed orders are $1$, $2$ and $4$, as predicted, and the cost required to reach an error of $10^{-6}$ differs by more than four orders of magnitude between the methods. Finally, we apply the Runge--Kutta method to the Lotka--Volterra predator--prey model and check the conservation of its first integral.

\bigskip
\noindent\textbf{Keywords:} ordinary differential equations, one-step methods, Runge--Kutta, order of convergence, global error.
\end{otherlanguage}
